\documentclass[../../../include/open-logic-section]{subfiles}

\begin{document}

\olfileid[pt]{sth}{story}{blv}

\olsection{Apêndice: Lei Básica V de Frege}

Na \olref[rus]{sec}, explicamos que Russell formulou o
paradoxo como um problema para o sistema que Frege delineou nos seus
\emph{Grundgesetze}. O sistema de Frege não incluía uma formulação direta
da compreensão ingênua. Logo, neste apêndice, explicaremos muito
sumariamente o que o sistema de Frege \emph{incluía}, como se relaciona com a compreensão ingênua e como se relaciona com o
paradoxo de Russell.

O sistema de Frege é de \emph{segunda ordem} e foi concebido para formular
a noção de \emph{extensão de um conceito}.\footnote{Estritamente
falando, Frege tenta formalizar uma noção mais geral: o
``valor de curso'' de uma função. As extensões de conceitos são um caso
especial da noção mais geral. Ver \citet[pp.\ 8--9]{Heck2012} para
pormenores.} Usando notação inspirada em Frege, escreveremos
$\fregeext{x}{F(x)}$ para \emph{a extensão do conceito~$F$}. Este
é um dispositivo que toma um \emph{predicado}, ``$F$'', e transforma-o num
\emph{termo} (de primeira ordem), ``$\fregeext{x}{F(x)}$''. Com este
dispositivo, Frege propôs a seguinte \emph{definição} de pertença:
\[
	a \in b =_\text{df} \exists G(b = \fregeext{x}{G(x)} \land Ga)
\]
grosso modo: $a \in b$ sse $a$ cai sob um conceito cuja extensão é
$b$. (Note que o quantificador ``$\exists G$'' é de segunda ordem.) Frege
também sustentava o seguinte princípio, conhecido como \emph{Lei básica V}: 
$$\fregeext{x}{F(x)} = \fregeext{x}{G(x)} \liff \forall x (Fx \liff Gx)$$
grosso modo: conceitos têm extensões idênticas sse são coextensivos. (De novo, tanto ``$F$'' como ``$G$'' estão em posição de predicado.) Um princípio simples liga agora a pertença à satisfação de propriedades:

\begin{lem}[em \emph{Grundgesetze}]\ollabel{lem:Fregeextensions}
$\forall F \forall a(a \in \fregeext{x}{F(x)} \liff Fa)$
\end{lem}

\begin{proof} 
Fixe $F$ e $a$. Ora $a \in \fregeext{x}{F(x)}$ sse $\exists G(\fregeext{x}{F(x)}
= \fregeext{x}{G(x)} \land Ga)$ (pela definição de pertença) sse
$\exists G(\forall x(Fx \liff Gx) \land Ga)$ (pela lei básica V) sse $Fa$
(por lógica elementar de segunda ordem).
\end{proof}

E disto resulta quase de imediato a compreensão ingênua:

\begin{lem}[em \emph{Grundgesetze}.]
$\forall F \exists s \forall a (a \in s \liff Fa)$
\end{lem}

\begin{proof}
Fixe $F$; agora \olref{lem:Fregeextensions} dá $\forall a (a \in
\fregeext{x}{F(x)} \liff Fa)$; logo $\exists s\forall a(a \in s \liff Fa)$ por
generalização existencial. O resultado segue-se, pois $F$ era
arbitrário.
\end{proof}

O paradoxo de Russell segue tomando $F$ dado por $\forall x(Fx \liff x \notin x)$. 

\end{document}
